\documentclass[11pt]{amsart}
\usepackage[margin=1.15in]{geometry}
\usepackage{amssymb,amsmath,amsthm,mathtools}
\usepackage[colorlinks=true,linkcolor=blue,urlcolor=blue,citecolor=blue]{hyperref}

\newtheorem{theorem}{Theorem}
\newtheorem{lemma}[theorem]{Lemma}
\newtheorem{proposition}[theorem]{Proposition}
\newtheorem{corollary}[theorem]{Corollary}
\theoremstyle{remark}
\newtheorem{remark}[theorem]{Remark}

% a step inside a proof: an underlined label with its claim in italics (stepclaim),
% followed by the proof of that claim, indented (stepproof)
\newenvironment{stepclaim}[1]{\par\smallskip\noindent\underline{\smash{Step #1}}:\enspace\itshape\ignorespaces}{\par}
\newenvironment{stepproof}{\begin{list}{}{\setlength{\leftmargin}{1.5em}\setlength{\rightmargin}{0pt}%
  \setlength{\topsep}{2pt}\setlength{\partopsep}{0pt}\setlength{\parsep}{0pt}\setlength{\itemsep}{0pt}%
  \setlength{\listparindent}{\parindent}}\item\relax}{\end{list}\smallskip}

\newcommand{\kk}{k}
\newcommand{\linf}{\ell^\infty}
\newcommand{\N}{\mathbb{N}}
\newcommand{\Z}{\mathbb{Z}}
\newcommand{\cL}{\mathcal{L}}
\newcommand{\cP}{\mathcal{P}}
\newcommand{\Alt}{\operatorname{Alt}}
\newcommand{\coeff}{\operatorname{coeff}}
\newcommand{\sh}{\operatorname{sh}}
\newcommand{\abs}[1]{\lvert #1\rvert}
\newcommand{\norm}[1]{\lVert #1\rVert}
\newcommand{\1}{\mathbf{1}}
\newcommand{\tail}[2]{#2|_{\ge #1}}

\title[Alternating precomposition is not analytic in characteristic 2]{Precomposition on continuous alternating maps\\ is not analytic in characteristic 2}
\author{Jack McCarthy}
\thanks{AI was used to assist in finding the result and in drafting this paper: the counterexample, its proof and the
Lean formalization were developed with Claude, ChatGPT and the AutoformBot plugin.}
% the date is printed under the author instead of in a first-page footnote
\makeatletter
\g@addto@macro\@setauthors{\vskip6pt{\centering\footnotesize October 2026\par}}
\makeatother

\begin{document}
\begin{abstract}
For normed spaces $E,E',F$ over a nontrivially normed field $\kk$ and a finite $\iota$, precomposition
$f\mapsto(m\mapsto m\circ(f,\dots,f))$ sends continuous linear maps $E\to E'$ to continuous linear maps between spaces of
continuous alternating maps. We show that it is \emph{not} analytic in general. If $\kk$ is a complete non-Archimedean field of
characteristic~$2$ that is not spherically complete, then for $E=E'=\linf$ and $\abs{\iota}=2$ there is a Banach space $F$
over $\kk$ for which this map is not analytic at $0$. The proof is a finite parity identity for determinants in
characteristic~$2$, glued along $\linf$. The result is formalized in Lean~4/Mathlib.
\end{abstract}
\maketitle

\section{Introduction}
Let $\kk$ be a nontrivially normed field. For normed spaces $E,E',F$ over $\kk$ and a finite $\iota$ we write $\cL(E,E')$
for the space of continuous linear maps with the operator norm, and $\Alt^\iota(E,F)$ for the normed space of continuous
alternating $\iota$-linear maps $E^\iota\to F$; we write $\Alt^2$ when
$\abs\iota=2$.

\subsection*{Analytic functors}
A map $f\colon X\to Z$ between normed spaces over $\kk$ is \textbf{analytic at} $x_0\in X$ if there are bounded $n$-linear maps
$p_n\colon X^n\to Z$ and a radius $r>0$ with $\sup_n\norm{p_n}r^n<\infty$ such that $f(x_0+h)=\sum_{n\ge0}p_n(h,\dots,h)$ whenever
$\norm h<r$. It is \textbf{analytic}, or $C^\omega$, if it is analytic at every point of $X$; this is the definition used in
Mathlib~\cite{docs}.

Lang~\cite[Ch.~III, \S4]{La} calls a functor $T$ on Banach spaces
\textbf{of class $\boldsymbol{C^p}$} if $f\mapsto T(f)$ is $C^p$. Similarly, we call a contravariant endofunctor $T$ on the category of normed spaces over $\kk$
\textbf{analytic} if for all $E,E'$ the map
\[
\cL(E,E')\to\cL\bigl(T(E'),T(E)\bigr),\qquad f\mapsto T(f),
\]
is analytic. For a normed space $F$ and a finite $\iota$, the assignment $E\mapsto\Alt^\iota(E,F)$,
$f\mapsto(m\mapsto m\circ(f,\dots,f))$, is a contravariant functor. The assignment is also functorial in $F$: a continuous linear map $g\colon F\to F'$ acts by $m\mapsto g\circ m$. This
action is linear in $g$, hence analytic, whereas the action of $f$ is not linear in $f$. So the bifunctor $\Alt^\iota$ is
analytic if and only if the functor $\Alt^\iota(-,F)$ is analytic for every $F$.

\subsection*{Motivation: alternating bundles in Mathlib}
In Mathlib, manifolds and vector bundles are modelled on normed spaces (of possibly infinite dimension) over an
arbitrary nontrivially normed field $\kk$, and their regularity is an element $n\in\N\cup\{\infty,\omega\}$, where
$C^\omega$ means analytic. Differential forms require the bundle of continuous alternating maps between two $C^n$ vector
bundles to be a $C^n$ vector bundle again (an instance of \texttt{ContMDiffVectorBundle}). Its transition functions are
those of the given bundles with $\Alt^\iota$ applied, so this comes down to the regularity of the functor $\Alt^\iota(-,F)$,
i.e.\ of the precomposition map
\begin{center}\texttt{ContinuousAlternatingMap.compContinuousLinearMapCLM}.\end{center}
Two cases have been settled, in pull requests to Mathlib.
\begin{itemize}
\item If $\abs\iota!$ is invertible in $\kk$ then $\Alt^\iota(-,F)$ is analytic, by polarisation. In Mathlib this is done in
characteristic~$0$ by Gou\"ezel~\cite{G1}, who uses it to construct the bundle for every $n$ in
characteristic~$0$~\cite{G2}; he remarks that a variant of the argument, not included there, covers finite-dimensional
spaces in positive characteristic, but not the general case.
\item Over every $\kk$, with no hypothesis on the characteristic or the dimension, $\Alt^\iota(-,F)$ is of class $C^n$ for every
$n\in\N\cup\{\infty\}$~\cite{PR}.
\end{itemize}
This leaves open the case where $\abs\iota!$ is not invertible in $\kk$ and the spaces are infinite-dimensional: is
$\Alt^\iota(-,F)$ analytic there? The following theorem provides a counterexample.

\begin{theorem}\label{thm:false}
Let $\kk$ be a complete non-Archimedean field of characteristic $2$ which is not spherically complete. There is a Banach
space $F$ over $\kk$ such that the functor $\Alt^2(-,F)$ is not analytic.
\end{theorem}

\section{Statement}\label{sec:statement}
\subsection*{Notation}
For a polynomial $f$ and a monomial $\mu$, $\coeff_\mu f$ is the coefficient of
$\mu$ in $f$. We write $c_0\subset\linf$ for the null sequences, $\1_S$ for the indicator of a set $S$, and $\1=\1_\N$.

Throughout, $\kk$ is a field which is complete for a nontrivial absolute value $\abs{\,\cdot\,}$ satisfying the
\textbf{ultrametric inequality} $\abs{x+y}\le\max(\abs x,\abs y)$; we call such a field \textbf{non-Archimedean}. The sup norms used below inherit the inequality. We assume that
\begin{itemize}
\item[(K1)] $\operatorname{char}\kk=2$;
\item[(K2)] $\kk$ is not spherically complete: there is a sequence of closed balls $B(c_n,r_n)$ such that
$\abs{c_{n+1}-c_n}\le r_n$, $r_0>r_1>\cdots$ and $\bigcap_n B(c_n,r_n)=\emptyset$.
\end{itemize}
The balls in (K2) are balls in $\kk$, and they are nested by the ultrametric inequality. Such fields exist; an explicit one is given in
Appendix~\ref{sec:field}. Let $E=\linf=\linf(\N,\kk)$ with the sup norm, a Banach
algebra under pointwise multiplication; write $m_a\in\cL(E,E)$ for multiplication by $a\in E$, so that
$\norm{m_a}=\norm a$, and $e_i$ for the unit vectors. For $x,y\in E$ let
\[
W(x,y):=x\otimes y+y\otimes x\in\linf(\N\times\N,\kk),\qquad W(x,y)_{pq}=x_py_q+x_qy_p ,
\]
and let $F\subseteq\linf(\N\times\N,\kk)$ be the closed linear span of all $W(x,y)$, a Banach space. Let
\[
A\colon \cL(E,E)\to\cL\bigl(\Alt^2(E,F),\Alt^2(E,F)\bigr),\qquad A(f)(m)=m\circ(f,f).
\]

\begin{theorem}\label{thm:main}
$A$ is not analytic at $0$. In particular $A$ is not $C^\omega$.
\end{theorem}

The map $A$ is the action on $\cL(E,E)$ of the functor $\Alt^2(-,F)$, so Theorem~\ref{thm:main} implies
Theorem~\ref{thm:false}. Sections~\ref{sec:step1}--\ref{sec:step4} prove Theorem~\ref{thm:main} in four steps.

\section{Step 1: analyticity forces a bilinear lift}\label{sec:step1}
Recall that $A\colon\cL(E,E)\to\cL(\Alt^2(E,F),\Alt^2(E,F))$ is given by $A(f)(m)=m\circ(f,f)$, so that $A(tf)=t^2A(f)$.
Steps~1 and~2 do not use \textup{(K2)}: until Step~3, $\kk$ may be any non-Archimedean field of characteristic~$2$.

\begin{lemma}\label{lem:lift}
Let $A\colon X\to Z$ be a map between normed spaces with $A(tf)=t^2A(f)$. If $A$ is analytic at $0$ then there is a bounded
bilinear $P\colon X\times X\to Z$ with $P(f,f)=A(f)$ for all $f$.
\end{lemma}
\begin{proof}
Let $\sum_np_n$ be a power series of $A$ at $0$, and fix $f\in X$. For scalars $t\in\kk$ with $\abs t$ small,
\[
t^2A(f)=A(tf)=\sum_nt^n\,p_n(f,\dots,f).
\]
Both sides are power series in the variable $t\in\kk$ with coefficients in $Z$. The coefficients of a convergent power
series in one variable are unique over any nontrivially normed field, so comparing the coefficients of $t^2$ gives
$p_2(f,f)=A(f)$. Take $P=p_2$.
\end{proof}

So assume that a bounded bilinear $P$ with $P(f,f)=A(f)$ exists. The map $W\colon E\times E\to F$ is bilinear, it is
alternating because $W(x,x)=2\,x\otimes x=0$, and it is continuous because $\norm{W(x,y)}\le\norm x\norm y$ by the
ultrametric inequality; so
$W\in\Alt^2(E,F)$ with $\norm W\le1$. Put, for $a,b,x,y\in E$,
\[
\Phi(a,b;x,y):=P(m_a,m_b)(W)(x,y)\in F .
\]
Then $\Phi$ is $4$-linear, and since $\norm{m_a}=\norm a$ and $\norm W\le1$ we have
\[
\norm{\Phi(a,b;x,y)}\le C\norm a\,\norm b\,\norm x\,\norm y,\qquad C:=\norm P.
\]
Here $C\ge1$, because $P(\mathrm{id},\mathrm{id})=A(\mathrm{id})$ is the identity map of $\Alt^2(E,F)\ne0$.
Moreover
\begin{equation}\label{eq:Phi}
\Phi(a,b;x,x)=0,\qquad \Phi(a,a;x,y)=W(ax,ay).
\end{equation}
Since $2=0$, the first identity gives the symmetry $\Phi(a,b;x,y)=\Phi(a,b;y,x)$. We will contradict the existence of such
a $\Phi$.

\section{Step 2: a finite parity identity}
Fix a finite set $I\subset\N$ of even size $2m$ and $a',b'\in E$ with $\norm{a'},\norm{b'}\le1$. Let
$X_i,Y_i$ ($i\in I$) be indeterminates. Expanding $\Phi$ multilinearly in $x=\sum_{i\in I}X_ie_i$ and $y=\sum_{i\in I}Y_ie_i$
defines, for any fixed $p,q\in I$, a polynomial
\[
G_{pq}:=\Phi(a'+x,\,b'+y;\,a'+x,\,b'+y)_{pq}\in\kk[X_i,Y_i]_{i\in I}
\]
in the polynomial ring in these $4m$ variables, which we abbreviate to $\kk[X,Y]$. The values of the $G_{pq}$ at points $x,y\in\kk^I$ (vectors supported on $I$) are the entries of an $I\times I$ matrix $G(x,y)$. We also put
\[
g_{pq}:=\Phi(x,y;x,y)_{pq}=\sum_{i,j,k,l\in I}\Phi(e_i,e_j;e_k,e_l)_{pq}\,X_iY_jX_kY_l\in\kk[X,Y] .
\]

\begin{proposition}\label{prop:symm}
The matrix $G=(G_{pq})_{p,q\in I}$ is symmetric with zero diagonal: $G_{pq}=G_{qp}$ and $G_{pp}=0$ in $\kk[X,Y]$.
\end{proposition}
\begin{proof}
Every element of $F$ is a symmetric matrix with zero diagonal, since $W(x,y)_{pq}=W(x,y)_{qp}$ and
$W(x,y)_{pp}=2x_py_p=0$, and both properties pass to the closed span. The coefficients of $G_{pq}$ are the $(p,q)$ entries
of values of $\Phi$, which lie in $F$.
\end{proof}

\begin{proposition}\label{prop:top}
The polynomial $g_{pq}$ is homogeneous of degree $2$ in the $X_i$ and of degree $2$ in the $Y_i$, and $G_{pq}-g_{pq}$ has
total degree $\le3$. In particular $G_{pq}$ has total degree $\le4$.
\end{proposition}
\begin{proof}
By $4$-linearity, $G_{pq}$ is a sum of $16$ terms, according to whether each of the four slots is filled by the constant
vector ($a'$ or $b'$) or by the variable one ($x$ or $y$), and the total degree of a term is the number of its variable
slots. So $g_{pq}$, which has two slots with $x$ and two with $y$, is exactly the degree $4$ part of $G_{pq}$, and
$G_{pq}-g_{pq}$ has total degree $\le3$.
\end{proof}

\begin{proposition}\label{lem:g}
For $p\neq q$ we have, in $\kk[X,Y]$,
\[
g_{pq}(Y,X)=g_{pq}(X,Y)+X_p^2Y_q^2+X_q^2Y_p^2\qquad\text{and}\qquad g_{pq}(X,X)=0 .
\]
\end{proposition}
\begin{proof}
The field $\kk$ is infinite, since it has a nontrivial absolute value, and a polynomial over an infinite field which
vanishes at every point is zero. So it suffices to check these identities at points $x,y\in\kk^I$. By the symmetry in the last two slots,
$g(y,x)=\Phi(y,x;y,x)=\Phi(y,x;x,y)$, so by bilinearity in the first two slots and \eqref{eq:Phi}
\begin{align*}
g(x,y)+g(y,x)&=\Phi(x+y,x+y;x,y)+\Phi(x,x;x,y)+\Phi(y,y;x,y)\\
&=W\bigl((x+y)x,(x+y)y\bigr)+W(x^2,xy)+W(xy,y^2)\\
&=\bigl[W(x^2,xy)+W(x^2,y^2)+W(xy,xy)+W(xy,y^2)\bigr]\\
&\qquad{}+W(x^2,xy)+W(xy,y^2)\\
&=W(x^2,y^2),
\end{align*}
where the last step uses $2=0$ and $W(xy,xy)=0$. This is the matrix whose $(p,q)$ entry is $x_p^2y_q^2+x_q^2y_p^2$. Finally
$g(x,x)=W(x^2,x^2)=0$.
\end{proof}

Let $\tau$ be the automorphism of $\kk[X,Y]$ swapping $X_i\leftrightarrow Y_i$, and $\delta$ the substitution $Y_i:=X_i$.

\begin{lemma}[balanced coefficients]\label{lem:bal}
Let $f\in\kk[X,Y]$ with $\tau f=f$. For every $S\subseteq I$ the coefficient of
$\prod_{i\in S}X_iY_i$ in $f$ equals the coefficient of $\prod_{i\in S}X_i^2$ in $\delta f$.
\end{lemma}
\begin{proof}
The coefficient of $\prod_{S}X_i^2$ in $\delta f$ is the sum of the coefficients in $f$ of all monomials $X^\alpha Y^\beta$ with
$\alpha+\beta=2\cdot\1_S$. The involution $(\alpha,\beta)\mapsto(\beta,\alpha)$ of this set of monomials preserves coefficients (as
$\tau f=f$), and its only fixed point is $\alpha=\beta=\1_S$; all other terms cancel in pairs. Thus the coefficient of
$\prod_SX_i^2$ in $\delta f$ is equal to the coefficient of $X^{\1_S}Y^{\1_S}=\prod_{i\in S}X_iY_i$ in $f$.
\end{proof}

A \textbf{perfect matching} of $I$ is a partition of $I$ into $m$ unordered pairs $\{p,q\}$, its edges; for example
$\{1,2,3,4\}$ has the three perfect matchings $12|34$, $13|24$ and $14|23$. Since $G_{pq}=G_{qp}$
(Proposition~\ref{prop:symm}), the product of the $G_{pq}$ over the edges of a perfect matching is well defined.

\begin{proposition}\label{prop:matching}
Let $M$ be a perfect matching of $I$ and $\nu:=\prod_{i\in I}X_iY_i$. Then
\[\coeff_\nu\prod_{\{p,q\}\in M}G_{pq}=1 .\]
\end{proposition}
\begin{proof}
Write $G_{pq}=g_{pq}+c_{pq}$, where $g_{pq}$ is homogeneous of total degree $4$ and $c_{pq}$ has total degree $\le3$
(Proposition~\ref{prop:top}). Expanding the product over the $m$ edges of $M$ gives $\prod_Mg_{pq}$ plus terms that contain
at least one factor $c_{pq}$, and each such term has total degree $\le3+4(m-1)=4m-1$. The monomial $\nu$ has total degree
$4m$, so only $\prod_Mg_{pq}$ contributes to its coefficient:
\[\coeff_\nu\prod_{M} G_{pq}=\coeff_\nu\prod_{M} g_{pq}.\]
Put $g'_{pq}:=g_{pq}+X_p^2Y_q^2$. By Proposition~\ref{lem:g},
\[\tau g'_{pq}=g'_{pq}\qquad\text{and}\qquad \delta g'_{pq}=X_p^2X_q^2 .\]
A product containing a factor $X_p^2Y_q^2$ has no squarefree monomial, so replacing every $g_{pq}$ by $g'_{pq}$ does
not change the coefficient of $\nu$. By Lemma~\ref{lem:bal}, applied to the $\tau$-invariant product $\prod_M g'_{pq}$,
this coefficient is the coefficient of $\prod_{i\in I}X_i^2$ in
\[\delta\prod_{M} g'_{pq}=\prod_{\{p,q\}\in M}X_p^2X_q^2=\prod_{i\in I}X_i^2,\]
which is exactly $1$.
\end{proof}

\begin{lemma}\label{lem:det}
Let $A$ be a symmetric matrix with zero diagonal over a commutative ring of characteristic $2$, of even size $2m$. Then
$\det A=\sum_M\prod_{\{p,q\}\in M}A_{pq}^2$, the sum being over perfect matchings; and the number of perfect matchings is
$(2m-1)!!$.
\end{lemma}
\begin{proof}
In characteristic $2$ we have $\det A=\sum_\sigma\prod_iA_{i\sigma(i)}$. As $A$ is symmetric, $\sigma$ and $\sigma^{-1}$ give the same product, so the permutations with
$\sigma\ne\sigma^{-1}$ cancel in pairs and only the involutions remain. An involution with a fixed point $i$ contributes the
factor $A_{ii}=0$. An involution without fixed points is a product of $m$ disjoint transpositions $(p\;q)$; the pairs
$\{p,q\}$ form a perfect matching $M$, every perfect matching arises exactly once in this way, and therefore
\[
\det A=\sum_M\prod_{\{p,q\}\in M}A_{pq}A_{qp}=\sum_M\prod_{\{p,q\}\in M}A_{pq}^2 .
\]
Finally, pairing the smallest index with one of the $2m-1$
others and continuing with the remaining indices shows that there are $(2m-1)(2m-3)\cdots3\cdot1=(2m-1)!!$
perfect matchings.
\end{proof}

\begin{lemma}[hypercube sum]\label{lem:cube}
Let $R$ be a commutative ring of characteristic $2$ and let $f\in R[Z_1,\dots,Z_N]$ have total degree $\le N$. Then
$\sum_{z\in\{0,1\}^N}f(z)$ is the coefficient of $Z_1Z_2\cdots Z_N$ in $f$.
\end{lemma}
\begin{proof}
By linearity we may take $f$ to be a monomial $Z^\alpha$. At a point $z\in\{0,1\}^N$ it takes the value $1$ if $z_i=1$ for all
$i$ in the support of $\alpha$, and $0$ otherwise; so the sum is $2^{N-\abs{\operatorname{supp}\alpha}}$, which vanishes in $R$
unless every variable occurs in $Z^\alpha$. The only monomial of degree $\le N$ in which all $N$ variables occur is
$Z_1\cdots Z_N$.
\end{proof}

\begin{corollary}\label{cor:point}
We have
\[\sum_{x,y\in\{0,1\}^I}\det G(x,y)=1,\]
and consequently there are $x,y\in\{0,1\}^I$ with $\abs{\det G(x,y)}\ge1$.
\end{corollary}
\begin{proof}
Each $G_{pq}$ has total degree $\le4$ (Proposition~\ref{prop:top}), so for a perfect matching $M$ the product $\prod_{\{p,q\}\in M}G_{pq}$ has total
degree $\le4m$. By Lemma~\ref{lem:cube} with $N=4m$, the number of variables of $\kk[X,Y]$, and by
Proposition~\ref{prop:matching},
\[
\sum_{x,y\in\{0,1\}^I}\ \prod_{\{p,q\}\in M}G_{pq}(x,y)=\coeff_\nu\prod_{\{p,q\}\in M}G_{pq}=1,
\qquad\text{where }\nu=\prod_{i\in I}X_iY_i .
\]
For each pair $x,y$ the
matrix $G(x,y)$ is symmetric with zero diagonal (Proposition~\ref{prop:symm}), so Lemma~\ref{lem:det} applies to it. In
characteristic $2$ we have the Frobenius endomorphism $a\mapsto a^2$ of $\kk$. Hence
\begin{align*}
\sum_{x,y\in\{0,1\}^I}\det G(x,y)
&=\sum_{x,y\in\{0,1\}^I}\ \sum_M\ \prod_{\{p,q\}\in M}G_{pq}(x,y)^2\\
&=\sum_M\Bigl(\,\sum_{x,y\in\{0,1\}^I}\ \prod_{\{p,q\}\in M}G_{pq}(x,y)\Bigr)^2=\sum_M1^2=(2m-1)!!=1,
\end{align*}
because $(2m-1)!!$ is odd. By the ultrametric inequality, $1\le\max_{x,y\in\{0,1\}^I}\abs{\det G(x,y)}$.
\end{proof}

\section{Step 3: locality on \texorpdfstring{$\linf$}{l-infinity}}
This step uses \textup{(K2)} and the completeness of $\kk$, which are standing hypotheses of \S\ref{sec:statement}, but not
the characteristic. Recall \textup{(K2)}: there are
centres $c_n\in\kk$ and radii $r_0>r_1>\cdots>0$ with $\abs{c_{n+1}-c_n}\le r_n$ such that no element of $\kk$ lies in all the
balls $B(c_n,r_n)=\{a\in\kk:\abs{a-c_n}\le r_n\}$. Recall also that $e_i\in\linf$ are the unit vectors. For $x\in\linf$ write $\tail Nx$ for $x$ with its coordinates below $N$ set to $0$.
\begin{proposition}\label{prop:nsc}
A bounded linear functional $\varphi$ on $\linf$ with $\varphi(e_i)=0$ for all $i$ is zero.
\end{proposition}
\begin{proof}
Suppose
$\varphi\ne0$. Since $r_{n+1}<r_n$ we have
$(r_{n+1}/r_n)\,\norm\varphi<\norm\varphi$, so by the definition of the operator norm there is $u_n\in\linf$ with
$\abs{\varphi(u_n)}>(r_{n+1}/r_n)\,\norm\varphi\,\norm{u_n}$. Put $s_n:=(c_{n+2}-c_{n+1})\,u_n/\varphi(u_n)$. Then
\[
\varphi(s_n)=c_{n+2}-c_{n+1}\qquad\text{and}\qquad
\norm{s_n}=\abs{c_{n+2}-c_{n+1}}\,\frac{\norm{u_n}}{\abs{\varphi(u_n)}}\le r_{n+1}\cdot\frac{r_n}{r_{n+1}\norm\varphi}=\frac{r_n}{\norm\varphi} .
\]
The difference $s_n-\tail n{s_n}$ is finitely supported, so by linearity $\varphi(s_n-\tail n{s_n})=0$; hence
$\varphi(\tail n{s_n})=c_{n+2}-c_{n+1}$. Moreover $\norm{\tail n{s_n}}\le\norm{s_n}\le r_n/\norm\varphi$. Define
$y:=\sum_k\tail k{s_k}$ coordinatewise; the sum is finite in every coordinate $i$ because $\tail k{s_k}$ vanishes at $i$ when
$k>i$. By the ultrametric inequality $\norm y\le\max_k\norm{s_k}\le r_0/\norm\varphi$, so $y\in\linf$. Likewise
$\norm{y-\sum_{k<n}\tail k{s_k}}=\norm{\sum_{k\ge n}\tail k{s_k}}\le\max_{k\ge n}\norm{s_k}\le r_n/\norm\varphi$. Hence, using the telescoping sum $\sum_{k<n}\varphi(\tail k{s_k})=c_{n+1}-c_1$,
\[
\abs{\varphi(y)+c_1-c_{n+1}}=\Bigl\lvert\varphi\Bigl(y-\sum_{k<n}\tail k{s_k}\Bigr)\Bigr\rvert
\le\norm\varphi\,\Bigl\lVert y-\sum_{k<n}\tail k{s_k}\Bigr\rVert\le\norm\varphi\cdot\frac{r_n}{\norm\varphi}=r_n .
\]
Put $a:=\varphi(y)+c_1\in\kk$. We have just shown that $\abs{a-c_{n+1}}\le r_n$, and $\abs{c_{n+1}-c_n}\le r_n$ by
\textup{(K2)}, so by the ultrametric inequality
$\abs{a-c_n}\le\max\bigl(\abs{a-c_{n+1}},\abs{c_{n+1}-c_n}\bigr)\le r_n$ for every $n$. Thus $a\in\bigcap_nB(c_n,r_n)$, which contradicts \textup{(K2)}.
\end{proof}
This is one direction of van der Put's theorem~\cite[Th\'eor\`eme~(3.2)]{vdP}: $\kk$ is spherically complete if and only if
$(\linf/c_0)'\ne0$. The next lemma, the description of the dual of $\linf$, is~\cite[Th\'eor\`eme~(4.1)]{vdP}; see
also~\cite[Ch.~4]{vR}. We include a proof.

\begin{lemma}\label{lem:dual}
Let $\varphi$ be a bounded linear functional on $\linf$. Then $\varphi(e_i)\to0$ as $i\to\infty$, and $\varphi(x)=\sum_ix_i\varphi(e_i)$ for every
$x\in\linf$. In particular $\abs{\varphi(u)}\le\norm u\,\sup_{i\ge N}\abs{\varphi(e_i)}$ whenever $u=\tail Nu$.
\end{lemma}
\begin{proof}
Suppose that $\varphi(e_i)\not\to0$. Then there are $\varepsilon_0>0$ and indices $i_0<i_1<\cdots$ with
$\abs{\varphi(e_{i_k})}\ge\varepsilon_0$. Let $T\colon\linf\to\linf$ be the operator with $T(x)(i_k)=x_k/\varphi(e_{i_k})$ and all other
coordinates $0$; it is linear and $\norm{T(x)}\le\norm x/\varepsilon_0$. So $\psi:=\varphi\circ T$ is a bounded functional, and
$\psi(e_k)=1$ for all $k$ because $T(e_k)=e_{i_k}/\varphi(e_{i_k})$. Let $\sh$ be the right shift, $\sh(x)(0)=0$ and $\sh(x)(i+1)=x(i)$. Then $\psi\circ(1-\sh)$ is a bounded functional with
\[
\bigl(\psi\circ(1-\sh)\bigr)(e_k)=\psi(e_k)-\psi(e_{k+1})=1-1=0\qquad\text{for every }k,
\]
so it vanishes by Proposition~\ref{prop:nsc}. But
$(1-\sh)\1=e_0$, so $0=\psi(e_0)=1$, a contradiction.

Hence $x_i\varphi(e_i)\to0$ for every $x\in\linf$. By the ultrametric inequality a difference of two partial sums of
$\sigma(x):=\sum_ix_i\varphi(e_i)$ is bounded by its largest term, so the partial sums form a Cauchy sequence, and since $\kk$
is complete the series converges, with $\abs{\sigma(x)}\le\norm x\,\sup_i\abs{\varphi(e_i)}\le\norm\varphi\norm x$. So $\sigma$ is a bounded functional with
$\sigma(e_j)=\varphi(e_j)$ for all $j$, so $\varphi=\sigma$ by Proposition~\ref{prop:nsc}.

For the last assertion let $u=\tail Nu$. Then $u_i=0$ for $i<N$, so
$\varphi(u)=\sum_{i\ge N}u_i\varphi(e_i)$. Every term of this series has absolute value at most
$\norm u\,\sup_{i\ge N}\abs{\varphi(e_i)}$; by the ultrametric inequality so does every partial sum, and hence so does the limit
$\varphi(u)$.
\end{proof}

\begin{proposition}[multilinear Phillips property]\label{prop:tails}
Let $\lambda\colon(\linf)^d\to\kk$ be a continuous multilinear map and $\varepsilon>0$. There is $N$ such that
$\abs{\lambda(u_1,\dots,u_d)}\le\varepsilon$ whenever $\norm{u_l}\le1$ for all $l$ and $u_l=\tail N{u_l}$ for some $l$.
\end{proposition}
We call such an $N$ a \textbf{threshold} for $\lambda$; every larger integer is then a threshold too.
\begin{proof}
Fix $\varepsilon$. By induction on $d$. For $d=1$ choose $N$ with $\abs{\lambda(e_i)}\le\varepsilon$ for $i\ge N$ and apply Lemma~\ref{lem:dual}.

Let $d\ge2$ and assume the statement for forms in fewer than $d$ variables. It suffices to show, for each $l$, that there
is $N_l$ such that $\abs{\lambda(u_1,\dots,u_d)}\le\varepsilon$ whenever $u_1,\dots,u_d$ lie in the unit ball of $\linf$ and
$u_l=\tail{N_l}{u_l}$; then $\max_lN_l$ is a threshold. Without loss of generality, assume $l=1$.
Suppose, for the sake of contradiction, that for every $N$ there are $u_1,\dots,u_d$ in the unit ball with
$u_1=\tail N{u_1}$ and $\abs{\lambda(u_1,\dots,u_d)}>\varepsilon$.

\begin{stepclaim}{1}
for every $N$ there are $i\ge N$ and $w_2,\dots,w_d$ in the unit ball with $\abs{\lambda(e_i,w_2,\dots,w_d)}>\varepsilon$.
\end{stepclaim}
\begin{stepproof}
For such a tuple $u_1,\dots,u_d$, Lemma~\ref{lem:dual} applied to the bounded
functional $\lambda(\cdot,u_2,\dots,u_d)$ gives
$\abs{\lambda(u_1,\dots,u_d)}\le\sup_{i\ge N}\abs{\lambda(e_i,u_2,\dots,u_d)}$, so one of the terms on the right is $>\varepsilon$.
\end{stepproof}

Define
\[
\nu_{l,i}(v_1,\dots,v_{d-1}):=\lambda(v_1,\dots,v_{l-1},e_i,v_l,\dots,v_{d-1})
\]
and $\mu_i:=\nu_{1,i}$. These are forms in $d-1$ variables, so the inductive hypothesis applies to them.

\begin{stepclaim}{2}
construct indices $j_1<j_2<\cdots$, integers $0=M_0<M_1<\cdots$ and vectors $z^{(s)}_2,\dots,z^{(s)}_d$ in the unit ball, supported in
$[M_{s-1},M_s)$ such that
\begin{equation}\label{eq:hump}
\abs{\mu_{j_s}(z^{(s)}_2,\dots,z^{(s)}_d)}=\abs{\lambda(e_{j_s},z^{(s)}_2,\dots,z^{(s)}_d)}>\varepsilon .
\end{equation}
\end{stepclaim}
\begin{stepproof}
Suppose that stages $1,\dots,s-1$ have been carried out. Let $K_s$ be a common threshold for the finitely many forms
$\nu_{l,i}$ with $2\le l\le d$ and $i<M_{s-1}$; for $s=1$ we put $K_1:=0$ and $j_0:=-1$. By Step~1 with
$N=\max(K_s,j_{s-1}+1)$ there are $j_s\ge K_s$, with $j_s>j_{s-1}$, and
$w_2,\dots,w_d$ in the unit ball with $\abs{\mu_{j_s}(w_2,\dots,w_d)}>\varepsilon$. Let $M_s>M_{s-1}$ be a threshold for $\mu_{j_s}$, and put
\[
z^{(s)}_l:=\tail{M_{s-1}}{w_l}-\tail{M_s}{w_l},\qquad\text{so that}\qquad w_l=h_l+z^{(s)}_l+t_l ,
\]
where $t_l:=\tail{M_s}{w_l}$ and $h_l=\sum_{i<M_{s-1}}c_{l,i}e_i$ with $\abs{c_{l,i}}\le1$. Expand $\mu_{j_s}(w_2,\dots,w_d)$ multilinearly into $3^{d-1}$ terms. A term containing some $t_l$ has absolute value
$\le\varepsilon$, because $M_s$ is a threshold for $\mu_{j_s}$. Consider a term containing no $t_l$ but
some $h_l$, say in slot $l$. By linearity in slot $l$, the term equals
\[
\sum_{i<M_{s-1}}c_{l,i}\,\lambda(e_{j_s},\dots,\underbrace{e_i}_{\text{slot }l},\dots)
=\sum_{i<M_{s-1}}c_{l,i}\,\nu_{l,i}(e_{j_s},\dots).
\]
The first entry is still
$e_{j_s}=\tail{K_s}{e_{j_s}}$, because $j_s\ge K_s$, and $K_s$ is a threshold for every $\nu_{l,i}$ with $i<M_{s-1}$. So every
summand has absolute value $\le\varepsilon$, and so does the term, by the ultrametric inequality. As
$\abs{\mu_{j_s}(w_2,\dots,w_d)}>\varepsilon$ while all these terms have absolute value $\le\varepsilon$, the ultrametric inequality shows
that the remaining term $\mu_{j_s}(z^{(s)}_2,\dots,z^{(s)}_d)$ satisfies~\eqref{eq:hump}.
\end{stepproof}

\noindent\underline{\smash{Step 3}}:\enspace because the supports $[M_{s-1},M_s)$ of the $z^{(s)}_l$ are pairwise disjoint, and each $z^{(s)}_l$ lies in the unit ball, the sum
$u_l:=\sum_sz^{(s)}_l$ is well defined with $\norm{u_l}\le1$. Let
$\chi:=\lambda(\cdot,u_2,\dots,u_d)$, a bounded functional. Fix $s$ and write $u_l=u_l^{<s}+z^{(s)}_l+u_l^{>s}$, where
$u_l^{<s}$ is supported in $[0,M_{s-1})$ and $u_l^{>s}$ in $[M_s,\infty)$. Expand
$\chi(e_{j_s})=\mu_{j_s}(u_2,\dots,u_d)$ multilinearly. Exactly as before, with $K_s$ and $M_s$ as in stage $s$, the terms
containing some $u_l^{>s}$ have absolute value $\le\varepsilon$, because $M_s$ is a threshold for $\mu_{j_s}$; and the terms containing
no $u_l^{>s}$ but some $u_l^{<s}$ have absolute value $\le\varepsilon$, because $K_s$ is a
threshold for every $\nu_{l,i}$ with $i<M_{s-1}$. The remaining term is
$\mu_{j_s}(z^{(s)}_2,\dots,z^{(s)}_d)$. By~\eqref{eq:hump} and the ultrametric inequality, $\abs{\chi(e_{j_s})}>\varepsilon$ for every
$s$. Since $j_s\to\infty$, this contradicts Lemma~\ref{lem:dual}.
\end{proof}
This is where (K2) is used. The statement fails on $c_0$, where $x\mapsto\sum_ix_i$ is a bounded functional taking the value
$1$ at every $e_i$, and it fails on $\linf$ over a spherically complete field, where the Hahn--Banach theorem~\cite{In}
gives a nonzero bounded functional vanishing on $c_0$.

\section{Step 4: gluing}\label{sec:step4}
Recall from Step~1 that $\Phi$ is $4$-linear with values in $F$ and that
\begin{equation}\label{eq:Phibound}
\norm{\Phi(a,b;x,y)}\le C\norm a\,\norm b\,\norm x\,\norm y\qquad\text{with }C\ge1 .
\end{equation}
We now combine Corollary~\ref{cor:point} (Step~2) with
Proposition~\ref{prop:tails} (Step~3).

Choose inductively finite sets $I_1<I_2<\cdots$ in $\N$ (every element of $I_t$ is smaller than every element of
$I_{t+1}$) with $\abs{I_t}=2t$ and vectors $a_t,b_t\in\{0,1\}^{I_t}$ as follows.
We write $a^{\le t}\in\linf$ for the vector that agrees with $a_s$ on $I_s$ for every $s\le t$ and is zero elsewhere, and define
$a^{<t}$, $b^{\le t}$ and $b^{<t}$ in the same way.
Given the earlier stages, apply Corollary~\ref{cor:point} with $a'=a^{<t}$ and $b'=b^{<t}$ to get $a_t,b_t$ with
$\abs{\det G_t}\ge1$, where $G_t=\Phi(a^{\le t},b^{\le t};a^{\le t},b^{\le t})|_{I_t\times I_t}$. Put $\varepsilon_t:=\tfrac12C^{-2t}$ and use
Proposition~\ref{prop:tails} (with $d=4$) for the finitely many entry forms $\Phi(\cdot)_{pq}$, $p,q\in I_t$, to get $N_t$ such that
these forms are at most $\varepsilon_t$ on unit arguments one of which is supported in $[N_t,\infty)$. All later sets $I_s$, $s>t$, are
chosen inside $[N_t,\infty)$.

Let $a,b\in\linf$ be defined by $a|_{I_t}=a_t$, $b|_{I_t}=b_t$ and zero elsewhere. Then
$a=\1_A$ and $b=\1_B$ are indicator functions of two subsets of $\N$, so $m_a,m_b$ are coordinate projections. Let
$\mathcal O:=\Phi(a,b;a,b)\in F$. Fix $t$ and split $a=a^{\le t}+a^{>t}$ and $b=b^{\le t}+b^{>t}$; the vectors
$a^{>t}$ and $b^{>t}$ are supported in the later sets $I_s$, $s>t$, hence in $[N_t,\infty)$. Expand
$\mathcal O$ multilinearly into $16$ terms. The term containing neither $a^{>t}$ nor $b^{>t}$ restricts to $G_t$ on
$I_t\times I_t$; each of the other $15$ terms has entries of absolute value $\le\varepsilon_t$ there, by the choice of $N_t$. So
$\norm{\mathcal O|_{I_t\times I_t}-G_t}\le\varepsilon_t$ by the ultrametric inequality. The entries of $\mathcal O$ and of $G_t$ have absolute value $\le C$ by~\eqref{eq:Phibound}, because $a$, $b$, $a^{\le t}$ and $b^{\le t}$ lie in the unit ball. Now expand
$\det\mathcal O|_{I_t\times I_t}=\det(G_t+\Delta_t)$, where $\Delta_t:=\mathcal O|_{I_t\times I_t}-G_t$, multilinearly in columns. One
term is $\det G_t$. Every other term contains a column of $\Delta_t$, whose entries have absolute value $\le\varepsilon_t$, so its
absolute value is at most $\varepsilon_tC^{2t-1}<1\le\abs{\det G_t}$. So the ultrametric inequality gives
\[
\abs{\det\mathcal O|_{I_t\times I_t}}=\abs{\det G_t}\ge1\qquad\text{for every }t .
\]
On the other hand, every element of $F$ is a uniform limit of matrices of finite rank, because each $W(x,y)$ has rank
$\le2$. So there is a matrix $R$ of some finite rank $r$ with $\norm{\mathcal O-R}\le\eta:=(2C)^{-1}$. Put $D:=\mathcal O-R$; its
entries have absolute value $\le\eta$, and those of $R=\mathcal O-D$ have absolute value $\le C$. Let $t>r$ and expand $\det(R+D)$ on
$I_t\times I_t$ multilinearly in columns, as before. A term with more than $r$ columns from $R$ vanishes, because these
columns are linearly dependent. Every other term has at least $2t-r$ columns from $D$, so its absolute value is at most
$C^r\eta^{2t-r}$, since $\eta\le1\le C$. By the ultrametric inequality,
$\abs{\det\mathcal O|_{I_t\times I_t}}\le C^r\eta^{2t-r}=2^{r-2t}C^{2(r-t)}<1$, which contradicts the previous display. This
proves Theorem~\ref{thm:main}. \qed

\appendix
\section{An explicit field}\label{sec:field}
The purpose of this appendix is to construct a field satisfying (K1) and (K2). Towards that end, let $k_0=\mathbb F_2$
(or any field of characteristic $2$). For
$j\ge0$ let $k_0((t^{1/2^j}))$ be the field of \emph{formal} Laurent series $\sum_{n\ge N}a_nt^{n/2^j}$ ($N\in\Z$, $a_n\in k_0$) in
the single variable $t^{1/2^j}$. These fields form an
increasing chain, and we let
\[
\cP:=\bigcup_{j\ge0}k_0((t^{1/2^j})),\qquad \abs f:=2^{-\operatorname{ord}f},
\]
be the field of \textbf{2-power Puiseux series}. Here $\operatorname{ord}f$ is the minimum $q$ for which $\coeff_{t^q}f\ne0$. We set $\abs0=0$ by definition.
Let $\kk$ be the completion of $\cP$ for this (non-Archimedean) absolute value. This is characteristic $2$ since
$k_0\hookrightarrow\cP\hookrightarrow\kk$.

\begin{proposition}\label{prop:field}
$\kk$ is not spherically complete.
\end{proposition}
\begin{proof}
We exhibit a sequence of balls as in (K2). Put
\[
\gamma_m:=1-2^{-(m+1)},\qquad
c_n:=\sum_{m=0}^{n}t^{\gamma_m}\in\cP,\qquad r_n:=2^{-\gamma_{n+1}} .
\]
\emph{The balls have the form required in \textup{(K2)}.} Since $c_{n+1}-c_n=t^{\gamma_{n+1}}$ we have
$\abs{c_{n+1}-c_n}=2^{-\gamma_{n+1}}=r_n$. Since the $\gamma_m$ are strictly increasing to $1$, the $r_n$ are strictly decreasing, and $r_n>2^{-1}=\tfrac12$.

\emph{The intersection is empty.} Suppose that $a\in\kk$ satisfies $\abs{a-c_n}\le r_n$ for all $n$. As $\cP$ is dense in $\kk$
there is $y\in\cP$ with $\abs{a-y}<\tfrac12$, and then for every $n$
\begin{equation}\label{eq:yball}
\abs{y-c_n}\le\max\bigl(\abs{y-a},\abs{a-c_n}\bigr)\le\max\bigl(\tfrac12,r_n\bigr)=r_n .
\end{equation}
Now $y\in k_0((t^{1/2^j}))$ for some $j$. The exponent
$\gamma_j=(2^{j+1}-1)/2^{j+1}$ has an odd numerator, therefore $\gamma_j$ is not a multiple of $1/2^j$ and so $\coeff_{t^{\gamma_j}}y=0$. For $n\ge j$ we have
$2^{-\gamma_j}>2^{-\gamma_{n+1}}=r_n$. Furthermore, the term $t^{\gamma_j}$ occurs in $c_n$, so
\[
\coeff_{t^{\gamma_j}}(y-c_n)=0-1\ne0,\qquad\text{hence}\qquad \operatorname{ord}(y-c_n)\le\gamma_j
\quad\text{and}\quad \abs{y-c_n}\ge2^{-\gamma_j}>r_n,
\]
which contradicts \eqref{eq:yball}.
\end{proof}

\section{Formalization}\label{sec:lean}
The Lean~4/Mathlib development is at
\begin{center}\url{https://github.com/Deicyde/alternating-analytic-counterexample}\end{center}
The library \texttt{AltCounter/} contains no \texttt{sorry}, and the only axioms it uses are the three standard ones
(\texttt{propext}, \texttt{Quot.sound} and \texttt{Classical.choice}). Its headline theorem assumes that Mathlib's
\begin{center}\texttt{ContinuousAlternatingMap.compContinuousLinearMapCLM}\end{center}
is $C^\omega$ for all nontrivially normed fields, normed spaces and finite index types, and derives \texttt{False}; a second
theorem does the same assuming this only for complete non-Archimedean fields of characteristic $2$, Banach spaces,
$E'=E$ and $\abs\iota=2$. Both are restated in a Mathlib-only file
\texttt{Challenge.lean} (with \texttt{sorry} in place of the proofs, as that format requires) and checked against
\texttt{Solution.lean} by the \texttt{leanprover/comparator} tool.

\begin{thebibliography}{99}
\bibitem{G1} S.~Gou\"ezel, \emph{In characteristic $0$, precomposition by a linear map is analytic on the space of
alternating maps}, pull request \#43338 to Mathlib,\\ {\small\url{https://github.com/leanprover-community/mathlib4/pull/43338}}.
\bibitem{G2} S.~Gou\"ezel, \emph{The bundle of continuous alternating maps is $C^n$}, draft pull request \#43234 to Mathlib,\\
{\small\url{https://github.com/leanprover-community/mathlib4/pull/43234}}.
\bibitem{In} A.~W.~Ingleton, \emph{The Hahn--Banach theorem for non-Archimedean-valued fields}, Proc.\ Cambridge Philos.\
Soc.\ \textbf{48} (1952), 41--45.
\bibitem{La} S.~Lang, \emph{Fundamentals of Differential Geometry}, Graduate Texts in Mathematics 191, Springer, 1999.
\bibitem{docs} The mathlib Community, \emph{Mathlib documentation}, \texttt{docs\#AnalyticAt},\\
{\small\url{https://leanprover-community.github.io/mathlib4_docs/find/?pattern=AnalyticAt#doc}}.
\bibitem{PR} J.~McCarthy, \emph{Smoothness reflects along a linear isometry with closed range}, draft pull request \#43548
to Mathlib,\\ {\small\url{https://github.com/leanprover-community/mathlib4/pull/43548}}.
\bibitem{vdP} M.~van der Put, \emph{Espaces de Banach non archim\'ediens}, Bull.\ Soc.\ Math.\ France \textbf{97} (1969), 309--320.
\bibitem{vR} A.~C.~M.~van Rooij, \emph{Non-Archimedean Functional Analysis}, Marcel Dekker, 1978.
\end{thebibliography}
\end{document}
